% !TeX root = ../main.tex
\section{Theoretical Foundations and Related Work}
\label{sec:estado_arte}

The proposed architecture converges three fundamental technological pillars: permissioned distributed ledgers with Byzantine Fault Tolerant (BFT) consensus, content-addressable distributed storage, and federated learning. Below, we synthesize their theoretical foundations and examine the critical limitations identified in the related literature.

\subsection{Permissioned Blockchain and IBFT 2.0 Consensus}
A blockchain is an immutable, distributed data structure where transactions are bundled into cryptographically linked blocks using collision-resistant hash functions (such as SHA-256)~\cite{nist2015sha, zheng2018blockchain}. Unlike public (permissionless) networks such as Bitcoin or Ethereum Mainnet, which suffer from high transaction latency and volatile monetary fees, industrial consortia predominantly adopt permissioned networks compatible with the Ethereum Virtual Machine (EVM)~\cite{buterin2014ethereum}, such as Hyperledger Besu~\cite{besu_docs}.

In these environments, compute-intensive (Proof-of-Work, PoW) or stake-based (Proof-of-Stake, PoS) mechanisms are replaced by Byzantine Fault Tolerant consensus algorithms. The \textbf{IBFT 2.0} (\textit{Istanbul Byzantine Fault Tolerant}) protocol~\cite{moniz2020istanbul} organizes consensus around a predefined set of validator nodes that alternate the proposer role in a round-robin schedule. A proposed block is appended to the ledger only after a three-phase voting process (pre-prepare, prepare, commit) signed by a two-thirds supermajority of validators:
\begin{equation}
n \ge 3f + 1
\end{equation}
where $n$ denotes the total number of validator nodes and $f$ represents the maximum number of tolerated faulty or Byzantine participants.

IBFT 2.0 provides a pivotal advantage over probabilistic consensus models: \textbf{immediate deterministic finality}. Because the chain does not experience probabilistic forks, once a training round commitment or model update is sealed within a block, it achieves immediate finality and permanent immutability on the ledger. Furthermore, while the transaction gas price is set to zero in private consortia to eliminate monetary friction, the computational metric of \textit{gas used}~\cite{ethereum_gas} remains an objective, deterministic measure of execution and storage overhead imposed on the EVM.

\subsection{Off-Chain Distributed Storage (IPFS)}
Directly persisting deep neural network parameters on a blockchain ledger is technically infeasible due to model artifact sizes (ranging from dozens to hundreds of megabytes), which would rapidly trigger severe EVM state bloat and demand prohibitive gas execution limits.

To overcome this bottleneck, we adopt the paradigm of \textbf{off-chain storage with on-chain commitment} utilizing the InterPlanetary File System (IPFS)~\cite{benet2014ipfs, daniel2022ipfs}. IPFS operates over a peer-to-peer (P2P) network based on content addressing and Merkle Directed Acyclic Graph (Merkle DAG) structures~\cite{daniel2022ipfs, trautwein2022design}: files are partitioned into immutable chunks and addressed through a unique Content Identifier (CID) derived directly from their cryptographic hash. The global-scale empirical and architectural evaluation conducted by Trautwein et al.~\cite{trautwein2022design} demonstrates that its block-exchange protocol (Bitswap) and distributed hash table (Kademlia DHT) achieve decentralized content routing with minimal overhead. By decoupling weight artifact storage to IPFS and recording only the content identifier (IPFS CIDv0, encapsulating the SHA-256 multihash of the serialized weights) within the smart contract~\cite{zhao2021privacy}, subsequent downloads can be validated against the immutable ledger. This provides data deduplication, formal integrity verification, and distributed access at negligible on-chain execution cost, provided that operational availability and persistence are backed by content pinning or dedicated IPFS nodes.

\subsection{Federated Learning and the FedAvg Algorithm}
Federated Learning (FL)~\cite{mcmahan2017federated} preserves data privacy by decoupling local model optimization from central parameter aggregation. Among federated consensus mechanisms, \textbf{FedAvg} (\textit{Federated Averaging}) serves as the foundational standard. At each communication round $t$, the aggregation coordinator broadcasts the current global weight vector $w^{(t)}$ to a selected set of $K$ edge clients. Each client trains the model on its private partition $\mathcal{D}_k$ by minimizing a local loss function, yielding updated local parameter weights $w_k^{(t)}$. The server then computes the consolidated global model as a sample-weighted average based on local sample counts $n_k$:
\begin{equation}
w^{(t+1)} = \sum_{k=1}^{K} \frac{n_k}{n} \, w_k^{(t)}, \quad \text{where } n = \sum_{k=1}^{K} n_k
\end{equation}

The \textbf{Flower} framework~\cite{beutel2020flower} offers a modular, efficient infrastructure to coordinate this federated lifecycle via bidirectional remote procedure calls over gRPC.

\begin{sidewaystable}
\caption{Qualitative and architectural comparison against relevant Blockchain-FL and robustness frameworks in the literature.}\label{tab:estado_arte}
\footnotesize
\setlength{\tabcolsep}{3.5pt}
\renewcommand{\arraystretch}{1.15}
\begin{tabular*}{\textheight}{@{\extracolsep\fill}lccccc}
\toprule
\textbf{Proposal / Framework} & \makecell{\textbf{Blockchain}\\\textbf{Consensus}} & \makecell{\textbf{Model}\\\textbf{Storage}} & \makecell{\textbf{Verification}\\\textbf{Protocol}} & \makecell{\textbf{Transactional}\\\textbf{Cost}} & \makecell{\textbf{Fault}\\\textbf{Tolerance}} \\
\midrule
\makecell[l]{Krum\\(Blanchard et al.)~\cite{blanchard2017byzantine}} & \makecell{No Blockchain\\(Centralized)} & \makecell{Central server\\(SPOF)} & \makecell{Geometric selection\\($L_2$)} & \makecell{None\\(No ledger)} & \makecell{Strict Byzantine\\($f < \frac{N-2}{2}$)} \\ \midrule
\makecell[l]{BlockFL\\(Kim et al.)~\cite{kim2019blockfl}} & \makecell{PoW\\(Slow / Inefficient)} & \makecell{On-Chain\\(Storage bottleneck)} & Unidirectional & \makecell{Volatile\\cryptocurrency} & \makecell{Probabilistic\\($< 50\%$)} \\ \midrule
\makecell[l]{DeepChain\\(Weng et al.)~\cite{weng2021deepchain}} & \makecell{PoW / PBFT\\with incentives} & \makecell{On-Chain\\(Encrypted gradients)} & \makecell{Homomorphic Enc.\\+ MACs} & \makecell{Severe (Massive\\on-chain ciphertext)} & \makecell{Formal Byzantine\\($f < N/3$)} \\ \midrule
\makecell[l]{AV-BFL\\(Pokhrel \& Choi)~\cite{pokhrel2020federated}} & \makecell{Adaptive PoW /\\DPoS} & \makecell{On-Chain\\at RSUs} & \makecell{Base station\\verification} & \makecell{Wireless\\retransmissions} & \makecell{Constrained by\\vehicular channel} \\ \midrule
\makecell[l]{H-BCFL\\(Chai et al.)~\cite{chai2021hierarchical}} & \makecell{Hierarchical DPoS\\at RSUs} & \makecell{On-Chain at\\base stations} & \makecell{Reputation and\\knowledge audit} & \makecell{Edge clustering\\overhead} & \makecell{Dependent on\\RSU reliability} \\ \midrule
\makecell[l]{B5G-FL\\(Lu et al.)~\cite{lu2020blockchain}} & \makecell{Permissioned\\(PBFT / BFT)} & \makecell{On-Chain at\\base stations} & \makecell{Base station\\consensus} & \makecell{DRL\\optimization} & \makecell{Formal Byzantine\\($f < N/3$)} \\ \midrule
\makecell[l]{Vehicular BFL\\(Pujol et al.)~\cite{pujol2024blockchain}} & \makecell{IBFT 2.0\\(Permissioned)} & \makecell{IoT alert\\logging} & \makecell{Traffic\\monitoring} & \makecell{Commercial\\gas-free} & \makecell{Formal Byzantine\\($N \ge 3f + 1$)} \\ \midrule
\makecell[l]{BCFL-Ref\\(Goh et al.)~\cite{goh2023blockchain}} & \makecell{PoW / PoA\\(Ethereum)} & \makecell{Hybrid Off-Chain\\(IPFS)} & \makecell{Unidirectional\\(No cross-verification)} & \makecell{Variable gas\\(Public/private net)} & \makecell{Formal Byzantine\\(Ethereum consensus)} \\ \midrule
\makecell[l]{\textbf{Our Architecture}} & \makecell{\textbf{QBFT / IBFT 2.0}\\\textbf{(Besu)}} & \makecell{\textbf{Content-addressable}\\\textbf{storage (IPFS)}} & \makecell{\textbf{Bidirectional Cross-Verif.}\\\textbf{+ Ex-Post Audit}} & \makecell{\textbf{Stable gas per}\\\textbf{transaction}} & \makecell{\textbf{Formal Byzantine}\\\textbf{($f=1 \to 2$)}} \\
\bottomrule
\end{tabular*}
\end{sidewaystable}

\subsection{Vulnerabilities in FL and Related Work}
Despite privacy advantages, conventional Federated Learning frameworks suffer from intrinsic vulnerabilities classified in recent security taxonomies and surveys~\cite{bouacida2021vulnerabilities, jere2021taxonomy, qu2022blockchain}:
\begin{enumerate}
    \item \textbf{Central server as a single point of failure and trust}: Existing frameworks offer no native mechanism to verify whether the central coordinator faithfully executes FedAvg or intentionally injects malicious biases into the global parameter vector $w^{(t+1)}$.
    \item \textbf{Poisoning attacks and lack of non-repudiation}: Compromised edge clients can execute model poisoning or inject stealthy backdoors~\cite{bagdasaryan2020backdoor}. In the absence of an immutable, non-repudiable log linking cryptographic identities to parameter digests, attributing malicious updates or auditing the history of compromised models becomes impossible.
\end{enumerate}

To address these vulnerabilities, researchers have increasingly investigated the hybridization of Blockchain and Federated Learning (BCFL)~\cite{qu2022blockchain}. However, designing practical architectures for Connected and Autonomous Vehicles (CAVs/IoV) involves four fundamental engineering trade-offs:

\textbf{1. On-Chain vs. Hybrid Off-Chain Storage and Computation:} Pioneer architectures such as BlockFL~\cite{kim2019blockfl} proposed exchanging full weight matrices directly over the blockchain ledger. Similarly, DeepChain~\cite{weng2021deepchain} integrated decentralized auditing and economic incentives with homomorphic encryption (HE) of gradients, incurring extreme computational and on-chain storage bottlenecks. Other designs, such as FLChain~\cite{bao2019flchain}, explored decentralized auditing using distributed Bloom filters, finding that serializing massive parameter tensors within blockchain transactions triggers acute EVM state bloat and degrades throughput. Consequently, the hybrid paradigm—leveraging content-addressable off-chain storage via IPFS anchored to on-chain cryptographic hashes~\cite{zhao2021privacy, goh2023blockchain}—represents an efficient approach to prevent ledger bloat while ensuring scalability in distributed deployments.

\textbf{2. Consensus Mechanisms and Latency in Wireless and Vehicular Networks:} In autonomous driving and edge computing (5G and beyond), communication latency and resource constraints are primary bottlenecks~\cite{pokhrel2020federated, lu2020blockchain}. Proof-of-Work (PoW) protocols introduce multi-minute confirmation delays and probabilistic forks that are intolerable in safety-critical vehicular systems~\cite{kim2019blockfl}. To optimize resource allocation and model convergence in cellular edge networks, Lu et al.~\cite{lu2020blockchain} formulated a base-station permissioned blockchain utilizing deep reinforcement learning (DRL) for dynamic bandwidth allocation, while Pokhrel and Choi~\cite{pokhrel2020federated} developed analytical models characterizing vehicular wireless channel dynamics on consensus delay. In the Internet of Vehicles (IoV), Chai et al.~\cite{chai2021hierarchical} proposed a hierarchical architecture offloading aggregation and DPoS consensus to Roadside Units (RSUs), Qi et al.~\cite{qi2021privacy} combined permissioned blockchain with differential privacy for traffic flow prediction, and Abou El Houda et al.~\cite{abouelhouda2024blockchain} devised collaborative intrusion detection schemes via BCFL. However, these frameworks depend heavily on the continuous availability and synchronization of fixed roadside infrastructure. Rather than relying on high-latency public networks or foregoing Byzantine fault tolerance, permissioned consortia running IBFT 2.0 / QBFT such as Hyperledger Besu~\cite{pujol2024blockchain} provide immediate deterministic finality and rigorous Byzantine resilience ($N \ge 3f + 1$) without monetary fee fluctuations.

\textbf{3. Cryptographic Integrity vs. Statistical Robustness:} Two complementary defensive tiers exist in federated learning security~\cite{bouacida2021vulnerabilities, jere2021taxonomy}. On the one hand, decentralized immutable ledgers provide formal guarantees regarding in-transit artifact integrity, sender authentication, and non-repudiation. On the other hand, if a participating client deliberately manipulates its local training data while legitimately signing and registering the resulting weights on-chain (\textit{data/model poisoning}), statistical aggregation defenses against anomalous gradients are essential. Literature has explored geometric filtering such as Krum~\cite{blanchard2017byzantine}, coordinate-wise trimming including Trimmed-Mean and Median~\cite{yin2018byzantine}, and hybrid defenses against complex collusions such as Bulyan~\cite{mhamdi2018hidden}. Similarly, FLTrust~\cite{cao2021fltrust} introduced trust bootstrapping by evaluating angular parameter similarity against a root dataset maintained by the coordinator. These two layers are complementary: on-chain non-repudiation unambiguously attributes each parameter update to a specific cryptographic identity, enabling post-facto forensic auditing and smart-contract-coordinated slashing once statistical anomalies are identified.

\textbf{4. Aggregator Accountability (Black-Box Aggregation vs. Ex-Post Auditing):} In verifiable aggregation literature without blockchain, pioneering works such as VerifyNet~\cite{xu2020verifynet} and VeriFL~\cite{guo2021verifl} address verifiable aggregation using homomorphic commitments, zero-knowledge proofs, and secret sharing schemes. Although these mechanisms formally ensure that the central coordinator computes aggregation faithfully without reconstructing raw gradients, they impose heavy computational overhead and synchronous communication delays on edge clients during every federated round. Conversely, most BCFL architectures either assume the central coordinator is trustworthy or attempt to execute aggregation directly within the EVM at prohibitive gas and computational costs. To the best of our knowledge, existing frameworks lack a decoupled, deterministic mechanism that—without penalizing operational cycle times ($<0.08\%$ overhead)—empowers any consortium participant or external auditor to reconstruct canonical FedAvg \textit{ex-post} from immutable Blockchain and IPFS commitments, mathematically verifying global model fidelity without compromising local data privacy.

To contextualize the contributions of this research with respect to prior work, Table~\ref{tab:estado_arte} presents a qualitative and architectural comparison. Unlike purely centralized statistical schemes (e.g., Krum), PoW-based architectures or massive on-chain storage paradigms (BlockFL, DeepChain), frameworks reliant on fixed RSU infrastructure, or unidirectional verification models lacking cross-validation (Goh et al.~\cite{goh2023blockchain}), our architecture implements bidirectional cross-verification paired with an independent ex-post audit protocol over a permissioned consortium blockchain with immediate finality and stable gas costs per transaction.
